\subsection{Modeled cash, reserves and the next plan}
The six-cycle direct-sale base produces 122 assurance, 121 planning and 119 dossier acceptances. These are simulated orders informed by the local execution bridge; they are not additional executed customer jobs. Net external worker collections are USD22,591.25; paid operating costs are USD11,621.14. Closing cash is USD12,970.11, including the original USD2,000 financing. Receivables of USD3,202.70 are excluded from spendable funds. After cumulative losses are recovered and the obligation reserve is removed, available earned operating surplus is USD8,776.09.

The prespecified next plan targets 20 orders per class, allowing up to two attempts each; acceptance is not guaranteed. Construction is USD439.69, operating working capital USD3,101.79, and one cycle of child maintenance USD23.35. Total precollection need is USD3,564.83. Full-surplus cash coverage is 2.46$\times$. However, its service cost exceeds the unchanged USD2,000 operating-spend ceiling: \textbf{cash coverage does not make this plan executable under the study's limits}.

A \textbf{post-result feasibility diagnostic}, without changing the ceilings, finds a largest balanced target of 12 orders per class, costing USD2,396.55 including construction and working capital. Cash coverage is 3.66$\times$. This is an illustrative resource-feasible target, not a newly prespecified success or guaranteed acceptance. The child still fails economic admission against ordinary caching.
\begin{table}[H]\small\centering\begin{tabular}{lrrrl}\toprule
Allocation & Earmarked USD & 20/class cash ratio & 12/class cash ratio & Decision \\\midrule
0\% & 0.00 & 0.00$\times$ & 0.00$\times$ & HOLD \\
25\% & 2,194.02 & 0.62$\times$ & 0.92$\times$ & HOLD \\
50\% & 4,388.04 & 1.23$\times$ & 1.83$\times$ & HOLD \\
75\% & 6,582.06 & 1.85$\times$ & 2.75$\times$ & HOLD \\
\bottomrule\end{tabular}\caption{Policy scenarios allocate an envelope, not permission to spend. All base policies retain the incumbent. The 20-order plan fails operating-spend admission; the smaller plan fits but its child has negative measured-compute value.}\end{table}
\begin{figure}[H]\centering\includegraphics[width=\linewidth]{figures/financing.pdf}
\caption{Modeled earned surplus after obligations, separate from seed cash and receivables. The two horizontal lines show full precollection plan needs. Crossing a cash line establishes neither plan feasibility nor positive investment value.}\end{figure}

\subsection{Admission, losing constructions and delayed benefits}
The measured child changes fresh-slice CPU by approximately $-3.027$ milliseconds of savings per joint request against conventional caching: it is slower. At the assumed compute rate, 100 uses have slightly negative compute value, while construction costs USD439.69 and three maintenance cycles add USD70.05. No human-attention or quality saving is invented to rescue the investment. The incumbent remains the recommended production method. A separately labeled compelled-research scenario can buy the child, but it lowers cash and does not establish payback.

The discovery-process pilot is also HOLD: its productive effect is unmeasured. In the provider-limit scenario, purchased capacity can increase assumed attempts when a provider ceiling actually binds. At 0\% allocation the model accepts 72/71/68 outputs by class; at 50\%, a USD350 capacity block increases them to 101/98/97. This is assumed vendor capacity scaling, not capability improvement. Benefits start the following cycle, expire after three cycles and require maintenance. Neither more cumulative output nor a higher resource ceiling is self-acceleration.

